% IMMUTABLE RESEARCH REPORT::FINAL-LINK2-LIT::2026-10-08_010300
\documentclass[11pt]{article}
\usepackage[T1]{fontenc}
\usepackage[margin=0.82in]{geometry}
\usepackage{amsmath,amssymb,amsthm,mathtools,booktabs,array,longtable,hyperref,xcolor}
\setlength{\emergencystretch}{3em}
\hypersetup{colorlinks=true,linkcolor=blue,citecolor=blue,urlcolor=blue}

\newtheorem{theorem}{Theorem}
\newtheorem{proposition}[theorem]{Proposition}
\newtheorem{definition}[theorem]{Definition}
\newcommand{\SAT}{\operatorname{SAT}}
\newcommand{\COVER}{\textsc{Cover-Shrink}}
\newcommand{\poly}{\operatorname{poly}}
\newcommand{\Open}{\textsc{Open}}
\newcommand{\NoGo}{\textsc{No-Go}}
\newcommand{\Positive}{\textsc{Positive}}

\title{FINAL-LINK2-LIT:\\Equality-Frontier Constructions and Exact No-Go Results\\for a Charged SAT Cover-Shrinker}
\author{RESEARCH agent report}
\date{8 October 2026}

\begin{document}
\maketitle

\begin{abstract}
This report updates FINAL-LINK-LIT with primary results on CSP
sparsification, streaming preprocessing, core reduction, polynomial
OR-kernels, truth-table reductions, parallel search-to-decision, and
fixed-solver unrolling.  The target is an exact constructor that, on every
ordinary CNF of encoded length $M$, emits at most $M^{1/3+o(1)}$ ordinary
SAT instances, each of length at most $M^{2/3+o(1)}$, while charging at most
$M^{1+o(1)}$ total discovery, output, replay, and local work.  The main new
positive result is that CSP core reduction is an efficient exact
preprocessing rule when a constraint language has domain redundancy.  The
main new negative result is sharper: for ordinary clause predicates,
nonredundancy characterizations show that retaining or reweighting a
subinstance cannot sparsify the violation objective below the dense scale;
one-pass exact SAT preprocessing likewise needs full dense-scale space.
Those theorems do not exclude derived clauses or a multi-query restriction
cover.  A recent circuit-lower-bound paper identifies deterministic
nonadaptive search-to-decision as an open bottleneck, closely matching the
WRAPPER's route-acquisition issue, but solving that bottleneck alone would
not shrink the encoded queries.  No primary result reviewed here constructs
the requested equality-frontier SAT cover.  The final link remains
\textbf{OPEN}, and the present work does not establish $P\ne NP$.
\end{abstract}

\section{The audited contract}

For an explicitly represented CNF $F$ of encoded length $M$, the desired
forward procedure must output ordinary explicit CNFs $F_1,\ldots,F_q$ and a
source-checkable combiner $\Phi$ such that
\begin{equation}
 \SAT(F)=\Phi\bigl(\SAT(F_1),\ldots,\SAT(F_q)\bigr),
\end{equation}
with
\begin{equation}
 q\le M^{1/3+o(1)},\qquad
 \max_i |F_i|\le M^{2/3+o(1)},\qquad
 W_{\rm total}\le M^{1+o(1)}.
\end{equation}
Here $W_{\rm total}$ includes discovery, certification, materialization,
replay, combination, and any preprocessing needed to formulate the child
queries.  A source pointer, program, or oracle gate that expands to the
entire parent is not an ordinary short CNF.  Such an object can be useful,
but it defines a succinct or typed oracle problem and requires a separate
return theorem to ordinary SAT.

The all-input quantifier and the simultaneous nature of (2) are essential.
A theorem that gives small children after exponential discovery, a small
number of full-size children, or a compact object whose expansion is large
does not meet the contract.  A randomized construction with nonzero failure
probability also does not by itself establish the deterministic exact
identity (1).

\section{The equality frontier is still a frontier, not a construction}

\begin{theorem}[Dell and van Melkebeek, Corollary 1]\label{thm:dvm}
Fix an integer $d\ge3$ and nonnegative reals $b,c$ with $b+c<d$.  Unless
$\mathrm{coNP}\subseteq\mathrm{NP/poly}$, there is no polynomial-time
reduction from $n$-variable $d$-SAT to an arbitrary problem that asks
$O(n^b)$ queries, each of bit length $O(n^c)$
\cite{dellvanmelkebeek2014}.
\end{theorem}

For a dense $n$-variable $3$-CNF, $M=n^{3+o(1)}$.  The target (2) becomes
\[
 q=n^{1+o(1)},\qquad |F_i|=n^{2+o(1)},
\]
so $(b,c)=(1,2)$ and $b+c=3=d$.  The target saturates the theorem's
inequality instead of violating it.  Dell and van Melkebeek describe their
tradeoff as tight at equality because the full source can always be sent as
one query.  That observation does not provide a balanced $n$-query,
$n^2$-bits-per-query construction.  The equality point remains available
but unfilled.

The sparse regime is more demanding.  If $M=n^{1+o(1)}$, a child of length
$M^{2/3+o(1)}$ cannot explicitly retain all $n$ variables.  An all-input
theorem therefore needs actual variable elimination, factorization, or a
semantic quotient.  A dense-only construction would not prove (2).

\section{New delta I: CSP sparsification separates two objectives}

Recent code and CSP sparsification results supply powerful near-linear
sparsifiers for some counting objectives.  They also give a direct reason
that ordinary SAT cannot use the most obvious version of this method.

\subsection{A constructive positive theorem}

\begin{theorem}[Khanna, Putterman, and Sudan, Theorem 1.1]
Let $C\subseteq\mathbb F_q^m$ be a dimension-$n$ linear code.  There is a
randomized algorithm running in
$\poly(n,m,\log q,\varepsilon^{-1})$ time that, with high probability,
computes a $(1\pm\varepsilon)$ code sparsifier of size
$\widetilde O(n/\varepsilon^2)$
\cite{khannaputtermansudan2024}.
\end{theorem}

Their CSP sparsifier preserves, for every assignment, the total weight of
satisfied constraints up to a multiplicative $(1\pm\varepsilon)$ factor.
For Boolean predicates they prove, among other classifications, an
efficient $\widetilde O_r(n^{r-1}/\varepsilon^2)$ bound when the predicate
does not have exactly one satisfying assignment, while a predicate with
exactly one satisfying assignment has no $o(n^r)$ sparsifier.  Their
Theorem 1.6 refines the exponent for ternary predicates through the largest
projection to an AND predicate; their Theorem 1.9 gives efficient
sparsification for predicates represented by low-degree polynomials over an
Abelian group.

This supplies a constructive pattern that a WRAPPER can reuse on promised
families: identify an importance or sensitivity measure, retain a weighted
representative subinstance, and stream the result.  It does not preserve
SAT exactly for general CNF.  Approximating the number of satisfied clauses
cannot distinguish $m$ satisfied clauses from $m-1$ at constant
$\varepsilon$.  Taking $\varepsilon=O(1/m)$ so that the gap becomes visible
changes the nominal sparsifier size from near-linear to at least quadratic
in $m$ and destroys the desired ledger.

\subsection{Why counting violations exposes the hard predicate}

One may instead count violated clauses.  A CNF is satisfiable exactly when
some assignment has violation count zero, so a multiplicative sparsifier of
violations would preserve the zero set.  For an $r$-clause, however, the
violation predicate is the unique falsifying assignment, equivalently a
signed AND or singleton predicate.  This is exactly the nonsparsifiable side
of the CSP classification.

\begin{theorem}[Brakensiek and Guruswami, Theorems 1.1 and 1.2]
For every nonlinear code $C\subseteq\{0,1\}^m$ and
$\varepsilon\in(0,1)$,
\[
 \operatorname{SPR}(C,\varepsilon)
   =O\!\left(\operatorname{NRD}(C)(\log m)^6/\varepsilon^2\right).
\]
For every nonempty proper predicate $R\subsetneq D^r$,
\[
 \operatorname{NRD}(\overline R,n)
 \le \operatorname{SPR}(R,n,\varepsilon)
 \le O\!\left(\operatorname{NRD}(\overline R,n)
                 (r\log n)^6/\varepsilon^2\right),
\]
where $\operatorname{NRD}$ is the maximum nonredundancy of the indicated
code or predicate family \cite{brakensiekguruswami2026}.
\end{theorem}

The same paper records
$\operatorname{NRD}(\operatorname{OR}_r,n)=\Theta(n^r)$.  Therefore the
complement predicate needed to preserve clause violations has full dense
nonredundancy.  In particular, a rule that only deletes or reweights source
clauses cannot turn every dense $r$-CNF into an exact subinstance of
$o(n^r)$ constraints while preserving all violation counts.  For $r=3$,
this blocks the hoped-for one-instance $n^2$ clause sparsifier inside that
restricted model.

This is an exact no-go for a specific construction class, not a no-go for
(1).  A multi-query cover may restrict variables differently in different
children.  A compiler may derive new clauses, introduce auxiliaries, or
preserve only the existence of a zero rather than every assignment's
violation value.  Those operations fall outside subinstance
sparsification.

Brakensiek and Guruswami also stress an algorithmic gap important to the
charge ledger.  Their general theorem is existential.  The naive
construction is polynomial in the code length, which can be exponential in
the number of CSP variables.  Obtaining a variable-scale efficient
construction remains open in general and is connected to CSP kernelization.
Thus even an existential small representative need not supply cheap
acquisition.

\section{New delta II: exact streaming preprocessing and core reduction}

\begin{theorem}[Sharma and Velusamy, Theorem 1.1]
For every fixed finite CSP language $\Gamma$ on a domain of size $q$ and
maximum arity $k$, $\SAT(\Gamma)$ has a deterministic one-pass streaming
algorithm using
$O_{q,k}(\operatorname{NRD}_n(\Gamma)\log n)$ space.  Every randomized
one-pass algorithm requires
$\Omega(\operatorname{NRD}_n(\Gamma))$ space
\cite{sharmavelusamy2026}.
\end{theorem}

For $k$-SAT, complete clause families give
$\operatorname{NRD}_n=\Theta(n^k)$, so an exact one-pass compact-state
preprocessor needs the full dense scale even with randomization.  The
positive upper bound maintains an equivalent nonredundant subinstance, but
its redundancy test enumerates all $q^n$ assignments because the theorem
places no restriction on time.  It therefore demonstrates that small state
and cheap discovery are independent obligations.

This streaming lower bound does not rule out all WRAPPERs.  A WRAPPER may
reread the source, use write-only aggregate output, or ask several oracle
queries.  The theorem does rule out treating an exact one-pass
delete-redundant-constraints procedure as the missing universal compiler.

A newer result identifies exactly when a domain-level positive rule exists.

\begin{theorem}[Brakensiek, Putterman, Sharma, and Velusamy, Theorems
1.1--1.3 and 2.1]
For every finite predicate $R$, deterministic one-pass
$\SAT(R)$ has space
$\widetilde O_R(\operatorname{NRD}_n(\operatorname{core}(R)))$, and every
randomized one-pass algorithm has matching
$\Omega_R(\operatorname{NRD}_n(\operatorname{core}(R)))$ space.  Replacing
each $R$-constraint by the corresponding $\operatorname{core}(R)$ constraint
on the same scope preserves satisfiability, the minimum number of
unsatisfied constraints, and the maximum number of satisfied constraints,
and can be performed in a stream.  Exact-CSP has, up to polylogarithmic
factors, space $n^d$, where $d$ is the degree of the core
\cite{brakensiekputtermansharmavelusamy2026}.
\end{theorem}

Core reduction is a genuine source-checkable exact contraction when a
constraint language contains redundant domain elements.  It is the
strongest new positive mechanism in this search.  Ordinary Boolean clause
languages are already core-like for this purpose, and signed $k$-clauses
retain full $\Theta(n^k)$ nonredundancy.  The construction therefore gives
a useful portfolio branch for structured CSPs, not the universal SAT
branch required by (2).

\section{New delta III: polynomial OR-kernels reveal a limited positive}

Fluschnik, Heeger, and Hermelin formalize a polynomial OR-kernel as a
nonadaptive family of oracle queries whose OR equals the source answer.
For parameterized Clique they give a simple closed-neighborhood kernel with
one query per vertex, each of size $O(\Delta^2)$, and substantially more
general polylogarithmic improvements.

\begin{theorem}[Fluschnik, Heeger, and Hermelin, Theorems 1 and 2]
For each real constant $c$, Clique admits polynomial OR-kernels with
$O(n/\log^c n)$ queries for the structural parameters studied in the paper.
Unless $\mathrm{NP}\subseteq\mathrm{coNP/poly}$, none of those
parameterizations admits a polynomial-size Turing kernel using
$O(n^{1-\varepsilon})$ queries for any constant $\varepsilon>0$
\cite{fluschnikheegerhermelin2022}.
\end{theorem}

The paper's generic lower-bound route says that a linear OR-composition plus
a sublinear-query polynomial Turing kernel gives an oracle
OR-distillation, and an oracle OR-distillation puts the language in
$\mathrm{coNP/poly}$.  This gives a precise lesson for parallel wrappers.
Local witness anchoring can reduce query count by arbitrary fixed powers of
$\log n$, but current generic techniques do not yield a polynomial power
saving.  Clique's local-neighborhood witness also has no direct analogue
that makes an arbitrary SAT assignment local.

The conditional lower bound does not exclude the target (2).  The target's
aggregate output is $M^{1+o(1)}$, and its query count is expressed in the
full input length $M$, rather than a parameter admitting the paper's linear
OR-composition.  It does show that merely placing several familiar kernels
in parallel is unlikely to create a hidden polynomial gain.

\section{New delta IV: truth-table normalization does not provide shrinkage}

\begin{theorem}[Buss and Hay]
Polynomial-time truth-table reducibility to SAT using Boolean circuits,
logspace truth-table reducibility to SAT using Boolean formulas, and
logspace Turing reducibility to SAT have the same power.  A constant number
of rounds of parallel SAT queries can be collapsed to one round
\cite{busshay1991}.
\end{theorem}

This theorem supports the double-wrapper architecture at the level of
adaptivity: a bounded number of query rounds can be normalized into a
nonadaptive truth-table family with an explicit Boolean combiner.  It gives
no bound that changes the number or encoded length of those queries.
Consequently it can flatten an already favorable wrapper, but it cannot
manufacture (2).

The modern batch-search literature makes the remaining distinction
especially clear.  Ren and Williams study time with bounded rounds of NP
queries and give a circuit-output algorithm based on the
Je\v{r}\'abek--Korten AVOID problem.

\begin{theorem}[Ren and Williams, Theorem 3.1, abbreviated]
Given the specified succinct truth-table circuit $G$ and length parameter
$T$, their AVOID algorithm either outputs a circuit of size
$O(|G|\log T)$ for the encoded function or outputs a nonimage, runs in
$\poly(T,|G|)$ time, and uses $O(n+\log(T/n))$ adaptive NP-query rounds with
witness length $n$ \cite{renwilliams2026}.
\end{theorem}

They identify the conversion of search-SAT to deterministic nonadaptive
decision-SAT queries as an open bottleneck.  Their stronger lower-bound
consequence assumes
$\operatorname{searchSAT}\in\mathrm{FP}^{\mathrm{NP}}_{\rm tt}$; the paper
notes that randomized efficient parallel search-SAT follows from isolation,
whereas the deterministic truth-table version is known under strong circuit
lower-bound assumptions and remains open at the desired level.

This is close to the WRAPPER's need to produce alternatives before learning
which one succeeds.  It is not the complete bridge.  Even a deterministic
nonadaptive search-to-decision reduction may issue full-size formulas.  It
could solve route acquisition while leaving the child-size inequality in
(2) untouched.

\section{Backdoor DNFs remain the closest exact cover syntax}

Ordyniak, Schidler, and Szeider define a $\mathcal C$-backdoor DNF as a
tautological DNF of partial assignments
$\tau_1,\ldots,\tau_q$ such that every residual $F[\tau_i]$ belongs to a
tractable class $\mathcal C$ \cite{ordyniakschidlerszeider2021}.  It gives
the exact identity
\[
 \SAT(F)=\bigvee_{i=1}^q\SAT(F[\tau_i]).
\]
The paper proves fixed-parameter detection in the number of terms for Horn,
AntiHorn, $2$CNF, and two unions of these classes.  Given a variable set
$B$, its displayed construction includes
$2^{3|B|+1}+3^{|B|}|F|^{O(1)}$ work.  This does not become near-linear when
the number of terms or backdoor variables grows as a positive power of
$M$.  Nor does the theorem guarantee that every CNF has a small backdoor DNF
or that every residual is smaller by a fixed power.

Backdoor depth is stronger than backdoor size on some structured instances.
The fixed-parameter approximation and bounded-depth linear-time evaluation
results for Horn, dual Horn, and Krom are valuable portfolio components
\cite{gaspersszeider2024}.  They remain parameter-promise results; they do
not bound backdoor depth on every formula.

Thus the literature still points to the following exact missing lemma.

\begin{definition}[Charged restriction cover]\label{def:cover}
A charged restriction cover for a CNF $F$ of length $M$ consists of partial
assignments $\tau_1,\ldots,\tau_q$ and a locally checkable coverage
certificate such that $\bigvee_i\tau_i\equiv1$,
\[
 q\le M^{1/3+o(1)},\qquad
 |F[\tau_i]|\le M^{2/3+o(1)},\qquad
 \sum_i|F[\tau_i]|\le M^{1+o(1)},
\]
and the terms, certificate, and explicit residuals are discovered and
replayed in $M^{1+o(1)}$ total work.
\end{definition}

A decision tree or disjoint cube partition can certify coverage
syntactically.  An arbitrary DNF does not: testing DNF tautology is
coNP-complete, and using a current-scale SAT call to establish coverage
reintroduces the circular query.  No reviewed source proves that every CNF
admits Definition~\ref{def:cover}.

\section{Representative sets and knowledge compilation do not change the delta}

Fomin, Lokshtanov, Panolan, and Saurabh compute a $q$-representative
subfamily of a family of independent $p$-sets in a rank-$(p+q)$ linear
matroid.  Its size is at most $\binom{p+q}{p}$ and their Theorem 1.1 gives
explicit-family cost
\[
 O\!\left(\binom{p+q}{p}t p^\omega
       +t\binom{p+q}{q}^{\omega-1}\right)
\]
field operations \cite{fominrepresentative2016}.  This preserves every
future disjoint extension, exactly the semantic property a valid merge
needs.  The theorem applies only after identifying a low-rank matroid whose
extension relation captures SAT.  General CNF satisfiability is not such a
relation under any construction supplied by the paper.  The explicit
family size $t$ in (5) is also charged.

Width-based knowledge compilation is similarly constructive on its promise
domain.  Amarilli, Monet, and Senellart compile a Boolean circuit with a
supplied width-$k$ tree decomposition into a d-SDNNF in
$O(|T|2^{(4+\varepsilon)k})$ time for every fixed
$\varepsilon>0$ \cite{amarillimonetsenellart2018}.  Canonical decision
diagrams and semantic hashing can merge equivalent residual states after
such a compilation.  Bova, Capelli, Mengel, and Slivovsky exhibit
polynomial-size CNFs requiring strongly exponential DNNF size
\cite{bovaetal2016}.  Their separation concerns full function
representation, so it does not rule out a decision-only multi-query cover,
but it rules out universal ``compile first, hash second'' as the cheap
acquisition mechanism.

Universal sets, covering codes, and split-and-list exhibit the same
boundary.  They provide exact coverage or sharing when the relevant
dimension is logarithmic or when exponential lists are allowed.  They do
not establish that an arbitrary CNF has only polynomially many relevant
assignment patterns.  Hash output length changes collision probability and
validation work, not the exponent of the number of semantic states.

\section{Fixed-solver specialization retains its exponent}

Suppose the contradiction hypothesis supplies a particular polynomial-time
SAT decider $A$ with running time
\[
 T_A(M)=M^{K+o(1)}.
\]

\begin{theorem}[Pippenger and Fischer, Theorem 4]
A computation of a one-dimensional-tape transducer taking $T$ steps can be
implemented by a combinational logic network of cost $O(T\log T)$ and depth
$O(T)$ \cite{pippengerfischer1979}.
\end{theorem}

Unrolling $A$ on length-$M$ inputs therefore produces a circuit of size
$M^{K+o(1)}$; a Tseitin encoding has the same polynomial scale.  This is a
useful upper bound and a legitimate specialization construction, but it
carries $K$.  It is not a lower bound against a different semantic
encoding.

The same issue appears in an explicit tableau split.  If a $T=M^K$ step
tableau is divided into $M^a$ blocks, each represented at scale $M^b$, then
the explicit-volume relation gives $a+b=K$.  If a recursive compiler reduces
a block exponent by a factor $b<1$, the exponent accumulated across levels
has the fixed-point form
\[
 \frac{a}{1-b}=\frac{K-b}{1-b}\ge K\qquad(K\ge1),
\]
before boundary-consistency and interface charges.  Equation (7) is an
accounting calculation, not a literature lower bound.  It shows why
ordinary subdivision of the known computation does not eliminate its
unknown polynomial exponent.

Instance compression under the qualitative assumption $P=NP$ has the same
fine-grained gap.  The assumption supplies some polynomial-time solver and
therefore some polynomial-time encoder, but it does not supply a uniform
bound independent of its exponent.  Calling that solver to decide which
restriction or representative is safe may cost $M^K$.  The desired proof
needs a structural acquisition theorem with the explicit bound in (2), not
only existence of a polynomial algorithm.

\section{Construction audit}

\renewcommand{\arraystretch}{1.2}
\begin{longtable}{@{}p{0.17\textwidth}p{0.22\textwidth}p{0.27\textwidth}p{0.26\textwidth}@{}}
\toprule
Mechanism & Exact positive content & Quantitative obstruction & Status for (1)--(2) \\
\midrule
\endhead
CSP satisfied-count sparsifier & Efficient $\widetilde O(n/\varepsilon^2)$ code sparsification and broad predicate classifications & Constant-error approximation cannot distinguish $m$ from $m-1$ satisfied clauses & Useful only on promised objectives; not exact SAT \\
Violation-count sparsifier & Zero would be preserved multiplicatively & Clause violation is a singleton/AND predicate with $\Theta(n^r)$ nonredundancy & Deletion/reweighting-only route blocked \\
CSP core reduction & Streaming, exact preservation of SAT, Min-CSP, and Max-CSP & Boolean clause language has no helpful domain redundancy and remains dense & Genuine positive for structured CSP, no general SAT shrinkage \\
One-pass exact preprocessing & Matching $\widetilde O(\operatorname{NRD})$ upper bound & $k$-SAT needs $\Omega(n^k)$ space; upper bound may use exhaustive time & Compact-state universal route blocked \\
Polynomial OR-kernel & Locally anchored Clique queries; polylogarithmic query savings & Power-saving query count conditionally excluded for studied parameters & Analogy only; no SAT locality theorem \\
Truth-table reduction & Bounded rounds can be flattened with an explicit combiner & No query-number or query-length improvement & Normalizes a good cover, does not create one \\
Parallel search-SAT & Randomized route through isolation; deterministic version studied explicitly & Deterministic $\mathrm{FP}^{\mathrm{NP}}_{\rm tt}$ search-SAT remains open at the needed level & Close acquisition clue, still no size reduction \\
Backdoor DNF & Exact heterogeneous restriction cover & No every-input small cover, residual-size bound, or near-linear detection & Closest syntax; missing theorem unchanged \\
Representative family & Exact preservation of future matroid extensions & Requires a low-rank SAT invariant; explicit family is charged & Possible implementation after a new invariant \\
Knowledge compilation & Exact merge and canonicalization under width promises & Exponential representations occur for general polynomial-size CNFs & Valid portfolio branch, not universal \\
Fixed-decider unrolling & Source-checkable circuit and CNF for a known computation & Size $M^{K+o(1)}$ retains the solver exponent & Cannot eliminate $K$ by itself \\
\bottomrule
\end{longtable}

\section{What changed and what did not}

FINAL-LINK2 adds three substantive facts to the previous audit.  First,
nonredundancy now gives a theorem-level obstruction to the most natural
clause-deletion and clause-reweighting attempts.  Second, core reduction is
a concrete exact positive construction that should be included in any
parallel portfolio, even though ordinary Boolean clauses do not benefit.
Third, deterministic nonadaptive search-to-decision is independently
recognized as an open bottleneck in current circuit-lower-bound work.  It
validates the architectural motivation for a forward-only wrapper but does
not supply encoded-size shrinkage.

The following facts did not change.  The target lies at, not below, the
Dell--van Melkebeek equality boundary.  OR-distillation lower bounds do not
exclude aggregate-linear multi-query output.  Backdoors, representative
sets, compilation, and hashing remain conditional on favorable source
structure.  Unrolling the assumed solver retains its exponent.  No reviewed
theorem simultaneously provides universal exactness, power-smaller
ordinary children, polynomially few children, and near-linear acquisition.

The strongest defensible remaining theorem is therefore
\[
 \boxed{\begin{minipage}{0.84\textwidth}
 For every explicit CNF $F$ of length $M$, construct a locally certified
 exact restriction or derived-instance cover with
 $q\le M^{1/3+o(1)}$, every ordinary child of length at most
 $M^{2/3+o(1)}$, and total discovery, materialization, replay, and local work
 $M^{1+o(1)}$.  The construction may use several portfolio branches, but it
 must prove that at least one branch succeeds on every input and must charge
 the work of all failed branches.
 \end{minipage}}
\]

Statement (8) is \Open.  It is not supplied by a known sparsification,
kernelization, backdoor, compilation, hashing, or self-reduction theorem.
If proved, recursive application would already give the recurrence
\[
 T(M)\le M^{1/3+o(1)}T(M^{2/3+o(1)})+M^{1+o(1)}
      =M^{1+o(1)},
\]
which explains why (8) is the central breakthrough rather than a routine
wrapper detail.  At present, the literature search does not complete a
defensible proof that $P\ne NP$.

\begin{thebibliography}{99}

\bibitem{dellvanmelkebeek2014}
Holger Dell and Dieter van Melkebeek,
\emph{Satisfiability Allows No Nontrivial Sparsification Unless the
Polynomial-Time Hierarchy Collapses},
Journal of the ACM 61(4), Article 23, 2014.
\url{https://eccc.weizmann.ac.il/report/2010/038/revision/1/download}

\bibitem{khannaputtermansudan2024}
Sanjeev Khanna, Aaron Putterman, and Madhu Sudan,
\emph{Efficient Algorithms and New Characterizations for CSP
Sparsification}, arXiv:2404.06327, 2024.
\url{https://arxiv.org/pdf/2404.06327}

\bibitem{brakensiekguruswami2026}
Joshua Brakensiek and Venkatesan Guruswami,
\emph{Redundancy Is All You Need (for CSP Sparsification)},
ECCC TR26-077, 2026; extended abstract in STOC 2025.
\url{https://eccc.weizmann.ac.il/report/2026/077/download}

\bibitem{sharmavelusamy2026}
Goutham Rajendran Sharma and Venkatesh Velusamy,
\emph{Characterizing Streaming Decidability of CSPs via Non-Redundancy},
arXiv:2604.21922, 2026.
\url{https://arxiv.org/pdf/2604.21922}

\bibitem{brakensiekputtermansharmavelusamy2026}
Joshua Brakensiek, Aaron Putterman, Goutham Rajendran Sharma, and
Venkatesh Velusamy,
\emph{Cores Characterize the One-Pass Streaming Complexity of CSPs},
arXiv:2610.07455, 2026.
\url{https://arxiv.org/pdf/2610.07455}

\bibitem{fluschnikheegerhermelin2022}
Till Fluschnik, Klaus Heeger, and Danny Hermelin,
\emph{Polynomial Turing Kernels for Clique with an Optimal Number of
Queries}, arXiv:2110.03279, 2021; conference version, 2022.
\url{https://arxiv.org/pdf/2110.03279}

\bibitem{busshay1991}
Samuel R. Buss and Louise Hay,
\emph{On Truth-Table Reducibility to SAT},
Information and Computation 91(1):86--102, 1991.
\url{https://doi.org/10.1016/0890-5401(91)90075-D}

\bibitem{renwilliams2026}
Hanlin Ren and R. Ryan Williams,
\emph{Near-Maximum Circuit Lower Bounds for Exponential Time with
Merlin-Arthur Queries}, ECCC TR26-118, 2026.
\url{https://eccc.weizmann.ac.il/report/2026/118/download}

\bibitem{ordyniakschidlerszeider2021}
Sebastian Ordyniak, Andre Schidler, and Stefan Szeider,
\emph{Backdoor DNFs}, IJCAI 2021, pages 1403--1409.
\url{https://www.ijcai.org/proceedings/2021/0194.pdf}

\bibitem{gaspersszeider2024}
Jan Dreier, Sebastian Ordyniak, and Stefan Szeider,
\emph{SAT Backdoors: Depth Beats Size},
Journal of Computer and System Sciences 142, 103520, 2024.
\url{https://doi.org/10.1016/j.jcss.2024.103520}

\bibitem{fominrepresentative2016}
Fedor V. Fomin, Daniel Lokshtanov, Fahad Panolan, and Saket Saurabh,
\emph{Efficient Computation of Representative Families with Applications in
Parameterized and Exact Algorithms},
Journal of the ACM 63(4), Article 29, 2016.
\url{https://arxiv.org/abs/1304.4626}

\bibitem{amarillimonetsenellart2018}
Antoine Amarilli, Mika\"el Monet, and Pierre Senellart,
\emph{Connecting Width and Structure in Knowledge Compilation},
ICDT 2018, Article 6.
\url{https://drops.dagstuhl.de/entities/document/10.4230/LIPIcs.ICDT.2018.6}

\bibitem{bovaetal2016}
Simone Bova, Florent Capelli, Stefan Mengel, and Friedrich Slivovsky,
\emph{A Strongly Exponential Separation of DNNFs from CNF Formulas},
ISAAC 2016.
\url{https://arxiv.org/abs/1411.1995}

\bibitem{pippengerfischer1979}
Nicholas Pippenger and Michael J. Fischer,
\emph{Relations Among Complexity Measures},
Journal of the ACM 26(2):361--381, 1979.
\url{https://doi.org/10.1145/322123.322138}

\end{thebibliography}

\end{document}
% END IMMUTABLE RESEARCH REPORT::FINAL-LINK2-LIT::2026-10-08_010300
